\documentclass[11pt]{article}
\usepackage[margin=1in]{geometry}
\usepackage{hyperref}
\usepackage{parskip}
\usepackage{enumitem}

\title{BMI Research --- Annotated Bibliography}
\author{Cole Zenk, Hayden, Kaden}
\date{\today}

\begin{document}
\maketitle

\begin{abstract}
This annotated bibliography surveys the current state of brain--machine
interface (BMI) research across three axes: the algorithms and digital
signal processing pipelines that turn raw neural recordings into usable
control signals, the technological and biological limitations that still
constrain implantable systems, and the clinical case studies that
demonstrate what BMIs can do today. Algorithmic sources cover spike
sorting, comparative neural decoding pipelines, and on-chip machine
learning for closed-loop prostheses. Limitations are examined through the
lenses of biocompatibility, device size and invasiveness, and signal
quality. Case studies include intracortical typing, cursor control via
Neuralink's N1 implant, and large-vocabulary speech decoding in patients
with ALS or paralysis. Together these sources frame the open question
driving this project: how machine learning and embedded hardware
co-design will shape the next generation of BMIs.
\end{abstract}

\section*{Research Questions}
\begin{itemize}
  \item How do BMIs work? What algorithms/techniques are used?
  \item What are the limitations of BMIs? Where have these shortcomings manifested in practice?
  \item How will ML impact BMIs?
\end{itemize}

\section*{Key Definitions}
\begin{description}
  \item[BMI] Brain machine interface.
\end{description}

% =====================================================================
\section{Notable Algorithms}
% =====================================================================

\subsection*{Spike sorting --- Lussac \cite{llobet2022lussac}}
This paper directly addresses a fundamental limitation in the spike-sorting step
of the neural signal processing pipeline: no single algorithm performs optimally
across all neuron types simultaneously. The authors introduce Lussac, a
Python-based post-processing tool that runs multiple spike-sorting passes ---
using tools like Kilosort and MountainSort with different filter bands and
detection thresholds --- then automatically merges the outputs using a
quality-guided cluster fusion algorithm. The quality metric combines cluster
contamination (estimated via refractory period violations) and firing frequency,
and merges are only accepted when they increase that score. Key DSP techniques
include bandpass filtering tailored to spike type (300--6000~Hz for simple spikes;
150--3000~Hz for complex spikes), cross-correlation-based spike alignment to
reduce detection jitter, and coincident spike deduplication across analyses.
Tested on cerebellar multielectrode recordings and the synthetic SpikeForest
dataset, Lussac yielded roughly a 50\% increase in detected Purkinje cell
clusters with both simple and complex spikes compared to either standalone
sorter. The paper is particularly relevant for understanding why spike sorting
is not a solved problem and how ensemble-based algorithmic approaches can
squeeze more signal out of noisy, low-SNR channels.

\subsection*{Comparative neural decoding \cite{shevchenko2024decoding}}
This paper conducts one of the most comprehensive head-to-head evaluations
of the full BCI decoding pipeline to date, testing 48 combinations of signal
processing, feature extraction, and classification across 14 subjects (672
total experiments). On the signal processing side, it compares artifact subspace
reconstruction (ASR), surface Laplacian filtering (SLF), and a normalization
baseline. Feature extraction methods include common spatial patterns (CSP),
independent component analysis (ICA), and short-time Fourier transform (STFT).
These are paired with four classifiers: SVM, LDA, CNN (EEGNet architecture), and
LSTM. The EEG dataset uses a 1--20~Hz bandpass filter and 200~ms rolling windows
with 160~ms overlap, optimized for decoding gait events from real-world walking.
Top-line findings: CNNs consistently outperformed all other classifiers (best F1
of $\sim$93\% with SLF+ICA); STFT feature extraction significantly degraded
performance across the board (up to 25\% drop); and ASR, despite being widely
used, often hurt accuracy or increased variance. Critically for embedded DSP
applications, the paper also benchmarks computational footprint --- CNN model size
was 0.03~MB versus 278.5~MB for SVM, and average inference time per segment was
under 1~ms for most feature extractors. This makes it a valuable reference for
evaluating algorithm trade-offs in real-time constrained systems.

\subsection*{On-chip ML for closed-loop prostheses \cite{zhu2021closedloop}}
This paper is the most hardware-oriented of the three, focusing on how ML-based
neural decoding algorithms are actually realized in silicon for real-time,
closed-loop prostheses. The review covers on-chip implementations across
epilepsy, Parkinson's disease, and other neurological applications, with
particular attention to application-specific integrated circuits (ASICs) and
algorithm-hardware co-design. The authors introduce an energy-area (E-A) figure
of merit to fairly compare across architectures accounting for channel count and
sampling rate. The core contribution is the Depth-Variant Tree Ensemble (DVTE),
a gradient-boosted decision tree classifier where trees of varying depths (1--4)
run in parallel, allowing shallow trees to output early predictions while deeper
trees correct errors. On a seizure detection task using intracranial EEG, DVTE
reduced processing latency by 2.5x compared to a conventional uniform-depth
ensemble with only a marginal drop in sensitivity ($<$3\%). A cost-aware learning
framework further penalizes power-hungry and high-latency features (e.g.,
low-frequency bandpower requiring long windows) during training, achieving a 3x
power reduction and 1.7x latency reduction at the same accuracy operating point.
The final DVTE classifier was implemented in TSMC 65~nm CMOS at $2.8~\mu W$
and $0.31~mm^2$, achieving 5.6~nJ/class energy efficiency --- over 6x better than
prior state-of-the-art. Directly relevant to the embedded DSP angle of this
project, as it demonstrates how algorithm design decisions (tree depth, feature
selection, quantization) have direct hardware consequences and must be
co-optimized from the start.

% =====================================================================
\section{Relevant Professions}
% =====================================================================

The development and improvement of these systems involve professionals such as
neurosurgeons, neuroscientists, and biomedical engineers, as evidenced by the
clinical implantation procedures, animal experiments, and material design
research described throughout the literature \cite{gao2025biocompat}.

% =====================================================================
\section{Technological Limitations}
% =====================================================================

\subsection*{Immune response and signal degradation \cite{gao2025biocompat}}
The article explains that the body's immune response leads to inflammation,
glial scar formation, and neuronal death, which degrade signal quality by
increasing the distance between electrodes and neurons and lowering the
signal-to-noise ratio. Evidence shows that microglia respond within 30 minutes
of implantation, and astrocytes form a glial scar over weeks, physically
isolating the device. Additional challenges include mechanical mismatch, where
rigid electrodes (in the GPa range) contrast sharply with soft brain tissue
($\sim$1~kPa), as well as micromotion that reduces cell viability and increases
tissue damage. The article also notes that oxidative stress from reactive oxygen
species can corrode electrode materials and contribute to device failure.

\subsection*{Size and invasiveness of current systems \cite{hoos2025columbia}}
Limitations of current BCI systems: existing devices require large implanted
canisters and invasive procedures such as removing parts of the skull or routing
wires through the body. Evidence shows that BISC addresses these issues by
reducing size to less than 1/1000th of conventional devices and using a thin
50~$\mu$m chip that can be inserted with minimal invasion, though challenges such
as long-term signal stability, tissue reactivity, and the need for continued
human testing still remain.

\subsection*{Signal quality, decoding, and regulatory hurdles \cite{kreg2025byfounders}}
Despite recent advancements, the article emphasizes significant technological
limitations. Non-invasive BCIs suffer from low signal quality due to interference
from the skull and scalp, while many systems still rely on outdated EEG
technology developed in 1924. Evidence shows that high-resolution alternatives
like MEG are bulky and expensive, and invasive systems face risks such as
infection, long-term complications, and declining signal stability due to
scarring and tissue movement. The article also highlights the difficulty of
decoding brain signals, comparing current systems to ``Morse code,'' as well as
challenges related to limited neural data, complex calibration requirements,
high manufacturing costs, and strict regulatory approval processes.

% =====================================================================
\section{Case Studies}
% =====================================================================

\subsection*{Case 1: Decoding attempted keyboard presses \cite{case_keyboard2026}}
The goal of this experiment was to decode attempted keyboard presses in patients
with paralysis by using brain activity to determine what keys they were trying
to press. This study used an intracortical Brain Computer Interface (BCI) typing
neuroprosthesis and tested it on two human participants. The first participant,
T17, was a 33 year old man with ALS. The second participant, T18, was a 48 year
old man with tetraplegia caused by a cervical spinal cord injury. Implanted
microelectrode arrays recorded neural activity from the precentral gyrus while
the participants attempted finger movements that matched QWERTY keyboard
presses. The system used 30 total keys, including the alphabet, space, and
punctuation. Different attempted finger flexion and extension movements were
mapped to different keys. A recurrent neural network decoded the neural activity
into character probabilities, and a 5-gram language model improved the final
sentence output. Before testing full sentence typing, the researchers showed
that individual attempted finger movements could be identified from brain
activity alone. Offline decoding accuracy for 30 isolated finger movements was
81.3\% for T17 and 95.2\% for T18. The main finding of the study was that both
participants were able to communicate through attempted typing, with the best
performance coming from T18. T18 reached a peak speed of 110 characters per
minute, which was about 22 words per minute, with a 1.6\% word error rate after
language model correction. T17 reached a peak of about 47 characters per minute.
The researchers concluded that this was the fastest reported hand-motor
intracortical BCI communication method at the time and that attempted finger
movements could be decoded accurately enough to support practical communication
in people with paralysis.

\subsection*{Case 2: Neuralink cursor control \cite{neuralink2024prime}}
The goal of this case study was to test whether a brain computer interface
could decode brain signals to allow for a paralyzed patient to control a
cursor. The participant was an adult with paralysis, and after implantation the
system was able to detect neural signals shortly after surgery. According to
the source, the participant then used the BCI system for digital tasks
including playing online chess and Sid Meier's Civilization VI. The system used
Neuralink's N1 Implant, which is an intracortical brain computer interface
designed to record neural activity through 1{,}024 electrodes distributed across
64 flexible threads. These signals were processed by electronics inside the
implant and then wirelessly transmitted to external software. That software
decoded the neural data into actions such as movement of a cursor on a computer
screen. The implant was powered by an internal battery and charged wirelessly,
allowing it to remain fully implanted under the scalp without external
connectors. For the implantation, researchers used functional MRI before
surgery to identify the area of the brain active during attempted hand and arm
movement. This was used to target the precentral gyrus for thread insertion.
The surgery involved exposing the target cortical region, using Neuralink's R1
surgical robot to insert the threads, and then mounting the implant into the
skull area before closing the scalp. The main finding from this case study is
that Neuralink was able to successfully implant a human BCI device, detect
usable neural signals, and use those signals to allow a participant with
paralysis to control digital applications through cursor-based interaction.

\subsection*{Case 3: Speech decoding in ALS \cite{willett2023speech}}
The goal of this experiment was to decode attempted speech in a patient who
could no longer speak clearly because of ALS. This study used an intracortical
speech brain computer interface and tested it on one participant, T12, who had
bulbar-onset ALS and could still make some limited orofacial movements and
vocalizations but could not produce intelligible speech. Microelectrode arrays
were implanted in speech-related areas of the brain, and the researchers
recorded neural activity while the participant attempted to speak phonemes,
words, and full sentences. The system worked by recording neural activity from
the implanted arrays and using a recurrent neural network to decode that
activity into phoneme probabilities over time. Those phoneme probabilities were
then combined with a language model to predict the most likely words and
sentences. Before testing full sentence decoding, the researchers first showed
that neural activity related to speech was strong enough to separate different
attempted speech movements. They were able to decode 33 orofacial movements
with 92\% accuracy, 39 phonemes with 62\% accuracy, and 50 words with 94\%
accuracy using simple classification methods. The main finding of the study was
that the participant was able to communicate through attempted speech at a much
faster rate than previous brain computer interface communication methods. The
participant reached 62 words per minute, with a 9.1\% word error rate on a
50-word vocabulary and a 23.8\% word error rate on a 125{,}000-word vocabulary.
According to the researchers, this was the first successful demonstration of
large-vocabulary speech decoding with this type of system and showed that
speech-related neural signals remained detailed enough years after paralysis to
support rapid communication.

\subsection*{Additional case: Handwriting BCI \cite{willett2021handwriting}}
A paralyzed patient was able to write using a BCI that interprets brain signals
during attempted handwriting motions, complementing the keyboard- and
speech-based decoding paradigms above.

\subsection*{Long-term recording and biocompatible electrodes \cite{gao2025biocompat}}
One study demonstrated long-term neural recording in epilepsy patients using
implanted platinum microwire electrodes for up to 15 months, showing clinical
feasibility and reliability. In animal studies, ultra flexible mesh electronics
injected into rodent brains showed minimal immune response and stable signal
recording over 4--12 weeks, while ECM-coated electrodes caused significantly
less brain damage and reduced scar formation three months after implantation.

\subsection*{Preclinical BISC testing \cite{hoos2025columbia}}
The article describes preclinical testing of the Biological Interface System to
Cortex (BISC), which demonstrated stable neural recording in motor and visual
cortices during surgical settings. Evidence from the study shows that the
implant has already been tested in real surgical environments and is moving
toward human use, with short-term intraoperative human studies currently
underway. Additionally, the system is being developed for treating
drug-resistant epilepsy through a National Institutes of Health--funded project
and has potential applications for conditions such as spinal cord injury, ALS,
stroke, and blindness by restoring motor, speech, and visual functions.

\subsection*{Real-world BCI applications survey \cite{kreg2025byfounders}}
The article describes several real-world applications of brain--computer
interfaces (BCIs) across medical and technological fields. For example,
speech-decoding systems developed by UCSF can convert neural activity into
real-time text or speech at speeds of up to 62 words per minute, approaching
normal conversation. Additionally, DARPA-funded projects have demonstrated
brain-controlled prosthetic limbs capable of multi-fingered movement, while
cochlear implants have already restored hearing for millions of patients.
Evidence from the article also includes the case of Ian Burkhart, whose implant
successfully restored function for years before being removed due to infection
risks, and Meta AI's Brain2Qwerty system, which achieved a 19\% character error
rate using non-invasive MEG and deep learning.

% =====================================================================
\section{Conclusion}
% =====================================================================

The literature surveyed here points to a field that has moved decisively
from proof-of-concept to early clinical reality, but one whose progress
is gated as much by hardware and biology as by algorithms. On the
algorithmic side, ensemble spike sorting~\cite{llobet2022lussac},
CNN-based decoders~\cite{shevchenko2024decoding}, and depth-variant tree
ensembles co-designed with silicon~\cite{zhu2021closedloop} all show
that meaningful gains in accuracy, latency, and power are still
achievable --- but typically at the cost of tighter coupling between
model and hardware. Case studies confirm the practical payoff: typing
at 22~words per minute~\cite{case_keyboard2026}, cursor control in
daily use~\cite{neuralink2024prime}, and 62~words per minute of decoded
speech~\cite{willett2023speech} were all unthinkable a decade ago.

At the same time, the limitations are sobering. Glial scarring,
mechanical mismatch, and oxidative degradation continue to erode signal
quality over months to years~\cite{gao2025biocompat}; non-invasive
modalities remain bandwidth-limited~\cite{kreg2025byfounders}; and
even state-of-the-art implants still trade off invasiveness against
longevity~\cite{hoos2025columbia}. The throughline across the sources
is that the next generation of BMIs will not be unlocked by any single
discipline: progress requires neurosurgeons, neuroscientists,
biomedical engineers, and ML researchers working against the same
constraints. This project's focus on the embedded DSP and ML layer sits
squarely in that intersection, and the works cited here form the
foundation for the technical questions to be pursued next.

% =====================================================================
\begin{thebibliography}{99}

\bibitem{llobet2022lussac}
Llobet, V., Wyngaard, A., \& Barbour, B. (2022).
\emph{Automatic post-processing and merging of multiple spike-sorting analyses with Lussac.}
bioRxiv. \url{https://doi.org/10.1101/2022.02.08.479192}

\bibitem{shevchenko2024decoding}
Shevchenko, O., Yeremeieva, S., \& Laschowski, B. (2024).
\emph{Comparative analysis of neural decoding algorithms for brain-machine interfaces.}
bioRxiv. \url{https://doi.org/10.1101/2024.12.05.627080}

\bibitem{zhu2021closedloop}
Zhu, B., Shin, U., \& Shoaran, M. (2021).
Closed-loop neural prostheses with on-chip intelligence: A review and a low-latency machine learning model for brain state detection.
\emph{IEEE Transactions on Biomedical Circuits and Systems, 15}(5), 877--897.
\url{https://doi.org/10.1109/TBCAS.2021.3112756}

\bibitem{gao2025biocompat}
Gao, W., et al. (2025, July 10).
\emph{Revolutionizing brain--computer interfaces: Overcoming biocompatibility challenges in implantable neural interfaces.}
Journal of Nanobiotechnology, U.S. National Library of Medicine.
\url{https://pmc.ncbi.nlm.nih.gov/articles/PMC12243264/}

\bibitem{kreg2025byfounders}
Kreg, M. (2025, February 18).
\emph{Brain-Computer Interfaces: Exploring the Last Frontier.}
byFounders VC.
\url{https://www.byfounders.vc/insights/brain-computer-interfaces}

\bibitem{hoos2025columbia}
Hoos, M. (2025, December 8).
\emph{Silicon chips on the brain: Researchers announce a new generation of brain-computer interface.}
Columbia Engineering.
\url{https://www.engineering.columbia.edu/about/news/silicon-chips-brain-researchers-announce-new-generation-brain-computer-interface}

\bibitem{case_keyboard2026}
\emph{Decoding attempted keyboard presses from intracortical neural activity in people with paralysis.}
Nature Neuroscience (2026).
\url{https://www.nature.com/articles/s41593-026-02218-y}

\bibitem{neuralink2024prime}
Neuralink. \emph{PRIME Study Progress Update.}
\url{https://neuralink.com/updates/prime-study-progress-update/}

\bibitem{willett2023speech}
Willett, F. R., et al. (2023).
A high-performance speech neuroprosthesis. \emph{Nature.}
\url{https://www.nature.com/articles/s41586-023-06377-x}

\bibitem{willett2021handwriting}
Willett, F. R., et al. (2021).
High-performance brain-to-text communication via handwriting.
\emph{Nature, 593,} 249--254.
\url{https://www.nature.com/articles/s41586-021-03506-2}

\end{thebibliography}

\end{document}
